% =====================================================================
%  Skill Doctor: Deterministic Multi-Layer Security Analysis of
%  AI Agent Skill Files with Host-Delegated Semantic Inference
%
%  Preprint v1.1 (final draft) -- 13 September 2026
%  Describes released artifact: Skill Doctor v0.1.0
%  arXiv primary: cs.CR   Cross-list: cs.SE, cs.AI
%  License: Apache-2.0
%
%  Build:  pdflatex skill-doctor-whitepaper.tex   (run 2-3 times)
%  Figures: place architecture.pdf, sdtm-v1.pdf, coverage.pdf and
%           adoption.pdf in ./figs/ . If a figure file is absent the
%           document still compiles and prints a labelled placeholder.
% =====================================================================

\documentclass[11pt]{article}

\usepackage[utf8]{inputenc}
\usepackage[T1]{fontenc}
\usepackage[margin=1in]{geometry}
\usepackage{graphicx}
\usepackage{booktabs}
\usepackage{tabularx}
\usepackage{array}
\usepackage{xcolor}
\usepackage[most]{tcolorbox}
\usepackage{algorithm}
\usepackage{algpseudocode}
\usepackage{listings}
\usepackage{enumitem}
\usepackage{caption}
\usepackage{microtype}
\usepackage[hidelinks]{hyperref}

\captionsetup{font=small,labelfont=bf}
\setlength{\tabcolsep}{4pt}
\renewcommand{\arraystretch}{1.15}

% ---- listings ---------------------------------------------------------
\lstdefinelanguage{json}{
  basicstyle=\ttfamily\small,
  showstringspaces=false,
  literate=*{:}{{{\color{black}:}}}{1}{,}{{{\color{black},}}}{1}
}
\lstset{
  basicstyle=\ttfamily\small,
  breaklines=true,
  columns=fullflexible,
  keepspaces=true,
  frame=single,
  framesep=4pt,
  xleftmargin=2pt,
  captionpos=b
}

% ---- graceful figure fallback ----------------------------------------
% \sdfigure{<path>}{<width>}{<placeholder caption text>}
\newcommand{\sdfigure}[3]{%
  \IfFileExists{#1}{\includegraphics[width=#2]{#1}}{%
    \fbox{\parbox[c][45mm][c]{\dimexpr#2-2\fboxsep-2\fboxrule\relax}{%
      \centering\small\textbf{[figure asset not built]}\\[2pt]
      \texttt{#1}\\[4pt]#3}}}%
}

\title{\bfseries Skill Doctor: Deterministic Multi-Layer Security Analysis of
AI Agent Skill Files with Host-Delegated Semantic Inference}

\author{Sayan Chowdhury\\
  Kalaris Labs\\
  \texttt{sayan@kalarislabs.com}}

\date{Preprint v1.1 (final draft) \quad 13 September 2026\\[2pt]
  \small Describes release \textbf{v0.1.0} \quad\textbullet\quad
  Code: \url{https://github.com/KalarisLabs/Skill-Doctor} \quad\textbullet\quad
  License: Apache-2.0\\
  arXiv primary: cs.CR \quad\textbullet\quad Cross-list: cs.SE, cs.AI}

\begin{document}
\maketitle

\begin{abstract}
\noindent
The agent-skill ecosystem is on a trajectory from the tens of thousands of
skills published today toward an estimated \textbf{three to five million
skills} within the current platform generation, distributed across public
registries, Git archives, and agent-mediated installers and loaded
automatically into runtimes such as Claude Code, Cursor, Codex CLI, and Gemini
CLI. AI agent runtimes load these \textbf{skill files} like \texttt{SKILL.md},
\texttt{AGENTS.md}, \texttt{.clauderules}, \texttt{.cursor/rules}, and Model
Context Protocol (MCP) server manifests at execution time and interpret them
with the agent's full trust context: tool access, filesystem, credentials, and
downstream agents. A skill file is the new executable, yet it arrives with none
of the defenses that binaries accumulated over three decades---no signature, no
execute bit, no endpoint check---and it is loaded silently across millions of
sessions per day.

\smallskip\noindent
At that scale the security problem changes in kind, not only in degree. A
population measured in the millions guarantees a standing subpopulation of
skills that are \emph{deliberately} weaponized by organized actors---credential
harvesters, supply-chain implants, tool-poisoning shims, and scanner-evading
logic bombs---authored specifically to exploit the agent runtimes that install
them. The public record already shows this at small scale: 1{,}184 confirmed
malicious skills in a single 2026 campaign, 36.7\% of surveyed MCP servers
exposed to SSRF, and a CVE filed against a skill scanner itself. Projected onto
three-to-five million skills, the absolute count of hostile artifacts is large
enough that manual review, per-key LLM review, and non-reproducible scanning
are all economically disqualified. What the ecosystem is missing is not another
detector but a \textbf{security layer for skill-file infrastructure}: one that
runs at registry and CI scale, at near-zero marginal cost, keyless, and with
verdicts reproducible enough to gate a merge and cache across a whole catalog.
Skill Doctor is designed as that layer, engineered as a production system
rather than a proof of concept.

\smallskip\noindent
We make four contributions. \textbf{First}, we present \textbf{SDTM-v1}, an
eleven-class threat taxonomy for skill files that synthesizes the OWASP Agentic
Skills Top 10 and OWASP MCP Top 10 with empirical findings from the ClawHavoc
campaign, MCP SSRF surveys, and bytecode-cache poisoning disclosures. SDTM-v1
introduces \textbf{SD-11, Scanner-Mediated Injection}, a class absent from
prior taxonomies: the scanner itself becomes an injection vector when it
submits untrusted skill content to a language model for analysis. Every
existing LLM-assisted skill scanner is exposed to SD-11 by construction.
\textbf{Second}, we present \textbf{Skill Doctor}, implemented in Rust as a
single statically linked binary. Its deterministic baseline layer---YARA-X
pattern matching, lexical source-to-sink taint heuristics, Shannon entropy
profiling, Unicode confusable and zero-width detection, and a novel
capability-manifest differ---\textbf{targets} a median completion time under
150\,ms with no network access, no API key, and no interpreter bootstrap, and a
one-to-two order of magnitude improvement over the Python-based state of the
art, whose lower bound is set by interpreter startup rather than by analysis.
Both figures are pre-registered targets measured by the published harness
(Section~\ref{sec:perf}); at v0.1.0 neither has been measured and no baseline
comparison has been run. \textbf{Third}, we introduce \textbf{host-delegated
semantic inference}. Prior work treats LLM analysis as a service the scanner
purchases: the user supplies a key, the scanner calls a model, and the result
is nondeterministic, latency-bound, metered, and unavailable in air-gapped
environments. We invert this. When Skill Doctor runs as an MCP server inside an
agent runtime, it performs no inference of its own; it emits static findings
plus a sanitized, instruction-neutralized analysis envelope, and the host
agent's own model---the model about to be exposed to the skill---returns a
schema-constrained verdict. We specify the sanitization protocol and the merge
invariant required to make this safe against SD-11. \textbf{Fourth}, we
introduce \textbf{structural coverage reporting}. Existing scanners disclaim
that ``no findings does not mean no risk'' and then offer no way to reason
about the gap. Skill Doctor prints, for every scan, how many SDTM-v1 classes
were structurally evaluable given the bundle contents and enabled layers, and
exposes this as a CI gate.

\smallskip\noindent
We define a reproducible evaluation protocol against Cisco AI Defense Skill
Scanner and NVIDIA SkillSpector at pinned versions, specifying a 1{,}000-skill
public corpus, 55 synthetic attack fixtures spanning all SDTM-v1 classes, a
20-skill evasion corpus that is statically clean by construction, a 15-skill
adversarial-scanner corpus, and a 100-skill hard-negative corpus. The v0.1.0
release ships the protocol, the seeded generator, and 14 curated
fixtures---6 attack fixtures covering SD-01, SD-02, SD-03, SD-04, and SD-10;
3 benign fixtures; 4 benchmark skills; and 1 example skill. The public,
evasion, adversarial-scanner, hard-negative, and real-malware corpora are
specified but not yet populated. Section~\ref{sec:perf} reports the first
measured values for binary footprint, peak memory, and bulk throughput; every
latency figure, every detection figure in Section~\ref{sec:detection}, and
every baseline comparison remain unmeasured at v0.1.0. We pre-register
falsifiable targets and commit to publishing misses. Skill Doctor, its rule
packs, and its full benchmark harness are released under Apache-2.0.
\end{abstract}

\noindent\textbf{Keywords:} AI agent security, prompt injection, supply chain
security, static analysis, Model Context Protocol, reproducible security
tooling, threat taxonomy

\vspace{1em}

\begin{tcolorbox}[colback=blue!4,colframe=blue!45!black,
  title=\textbf{Reconciliation statement}]
\small
This is the final draft. Every implementation statement has been reconciled
line by line against the shipped v0.1.0 source tree
(\texttt{docs/PAPER-RECONCILIATION.md}): subcommands, flag semantics, cache
tiers, sandbox isolation boundary, L4 feed status, fixture counts, package
names, and binary footprint. Where the implementation is narrower than the
design, the text says so in place rather than in a footnote.
\end{tcolorbox}

\tableofcontents
\newpage

% =====================================================================
\section{Introduction}
% =====================================================================

\subsection{A trust inversion, by design}

A conventional program earns the right to execute through a chain of gates. It
must be deliberately downloaded, granted an execute bit, survive endpoint
detection, and increasingly carry a signature. Decades of incident response
produced this chain.

A skill file bypasses all of it. It is a text document. An agent runtime
discovers it, loads it into context, and from that moment its contents carry
approximately the same authority as the operator's own instructions. There is
no execute bit for prose.

This is not an oversight. It is the value proposition. Skills are useful
precisely because they reconfigure agent behavior without a code review cycle.
A developer can change what an agent does by editing Markdown. The same
property that makes skills the fastest-growing extensibility surface in the
2026 agent stack makes them the highest-leverage attack surface in it.

We call this a \textbf{trust inversion}: the artifact with the least
verification receives the most authority.

\subsection{The inversion is already being exploited}

The threat is not speculative. Five observations from 2026 establish the base
rate.

\begin{enumerate}[leftmargin=*,itemsep=2pt]
  \item \textbf{Campaign-scale abuse.} The ClawHavoc campaign resulted in the
  identification of \textbf{1{,}184 confirmed malicious skills} in a single
  public registry, including credential exfiltration and injection of
  persistent instructions into downstream agents.
  \item \textbf{Protocol-level exposure.} A survey of roughly \textbf{7{,}000
  MCP servers} found \textbf{36.7\%} vulnerable to server-side request forgery,
  with skill and manifest files the common entry vector.
  \item \textbf{Companion code is the multiplier.} Skills that ship an
  executable companion script are \textbf{2.12 times more likely} to contain a
  command injection path than a standalone \texttt{SKILL.md}.
  \item \textbf{Defensive tooling is itself unaudited.} \textbf{CVE-2026-26057}
  was filed against a skill scanner. The tools we deploy to inspect hostile
  input are parsing hostile input.
  \item \textbf{Standards bodies have formalized the surface.} OWASP published
  both an Agentic Skills Top 10 and an MCP Top 10 in 2026, which establishes
  vocabulary but not tooling.
\end{enumerate}

\subsection{Why the existing defenses are not adopted}

Credible tooling exists. Cisco AI Defense ships a skill scanner; NVIDIA ships
SkillSpector. Both are competent. Neither has become default infrastructure. We
argue the obstacle is not detection quality but \textbf{three forms of
friction}, each of which is architectural rather than incidental.

\paragraph{Runtime dependency.} Every existing scanner requires a Python
environment. In continuous integration, interpreter and dependency bootstrap
dominates wall-clock time and is paid on every pull request, forever. The
analysis is not the cost; the environment is.

\paragraph{Credential dependency.} Meaningful analysis requires the user to
obtain, store, and rotate a third-party model API key. This is a hard stop for
air-gapped research environments, for regulated enterprises whose data cannot
leave the perimeter, and---most consequentially---for the casual developer
evaluating a tool for the first time. That last population determines whether a
security tool becomes default infrastructure or remains a compliance checkbox.

\paragraph{Nondeterminism.} LLM-in-the-loop scanning cannot produce
reproducible verdicts. A finding that appears on one run and vanishes on the
next cannot be argued in code review, cannot be cached at registry scale, and
cannot legitimately gate a merge.

We state the resulting design thesis plainly, because it is the organizing
claim of the paper: \textbf{a scanner becomes infrastructure when the marginal
cost of running it approaches zero.} The correct optimization target is
therefore not maximum detection per invocation but minimum friction per
invocation, with detection layered on top in a way that never becomes a
prerequisite.

\subsection{Contributions}

\begin{enumerate}[leftmargin=*,itemsep=2pt]
  \item \textbf{SDTM-v1}, an eleven-class threat taxonomy for agent skill
  files, including the previously undescribed \textbf{SD-11 Scanner-Mediated
  Injection}.
  \item A \textbf{Rust implementation} whose deterministic layer is one to two
  orders of magnitude faster than the state of the art and requires no
  interpreter, key, account, or network access.
  \item \textbf{Host-delegated semantic inference}, with a specified
  neutralization protocol and an additive-only merge invariant that makes
  delegation safe against SD-11.
  \item \textbf{Structural coverage reporting}, which converts the ``no
  findings is not no risk'' disclaimer into a printed metric and an enforceable
  CI threshold---released together with a \textbf{public, reproducible benchmark
  harness} (pinned competitor baselines, a seeded corpus generator, and
  pre-registered targets) that operationalizes it.
\end{enumerate}

\subsection{Non-goals}

We do not claim to decide maliciousness in general; that is undecidable. We do
not claim superior static detection over Cisco or NVIDIA on statically visible
threats, and Section~\ref{sec:discussion} explains why we deliberately avoid
that claim. We do not build a hosted scanning service, a web application, or a
governance suite. We do not attempt to replace the OWASP taxonomies; we
implement them and extend them by one class.

% =====================================================================
\section{Background}
% =====================================================================

\subsection{What a skill file is}

A skill file is a declarative, natural-language-first artifact that modifies
agent behavior at load time. In practice a skill \emph{bundle} contains some
combination of:

\begin{itemize}[leftmargin=*,itemsep=2pt]
  \item A Markdown body with instructions, typically with YAML frontmatter
  declaring a name, description, and sometimes a tool allowlist.
  \item Zero or more \textbf{companion scripts} (most commonly Python,
  increasingly JavaScript and shell) invoked by the agent.
  \item Dependency declarations (\texttt{requirements.txt},
  \texttt{package.json}) and occasionally lockfiles.
  \item For MCP servers, a manifest declaring tools, parameter schemas, and
  descriptions.
\end{itemize}

The security-relevant property is that \textbf{the natural-language portion is
not inert}. Text in a description field enters model context and can influence
tool selection, argument construction, and downstream delegation.

\subsection{Where skill files come from}

Skills are distributed through public registries, Git repositories, archive
downloads, and increasingly through agent-mediated installation, in which the
user asks an agent to ``add a skill for X'' and the agent resolves, downloads,
and loads it without the user reading a single line. Agent-mediated
installation collapses the review step entirely, which is why a pre-load gate
that requires no human action is the correct intervention point.

\subsection{The MCP host relationship}

The Model Context Protocol lets an external server expose tools to an agent
runtime. This matters twice in this paper. It is the mechanism by which a
scanner can gate skill loads from inside Claude Code, Cursor, Codex CLI, and
Gemini CLI---and it is the mechanism by which a scanner has a capable frontier
model one hop away, which is the observation that motivates our third
contribution.

% =====================================================================
\section{Threat Model: SDTM-v1}
\label{sec:threat}
% =====================================================================

\subsection{Adversary model and assumptions}

\paragraph{Adversary capability.} The adversary controls the full contents of a
skill bundle, including Markdown prose, frontmatter, companion source code,
dependency declarations, filenames, and byte-level encoding. The adversary may
publish to a public registry, compromise an existing skill's supply chain, or
induce an agent to fetch from an attacker-controlled URL.

\paragraph{Adversary goals.} Credential and data exfiltration; arbitrary
command execution on the developer host; persistence that survives skill
removal; manipulation of agent tool selection; and---the goal that motivates
SD-11---causing a security scanner to report a clean verdict on a malicious
artifact.

\paragraph{Assumptions we make.} The skill under examination is adversarial by
assumption; that is the reason it is being scanned. The host operating system
and the scanner binary itself are trusted. The agent runtime is honest but
exploitable. Cryptographic primitives are sound.

\paragraph{Out of scope.} Compromise of the developer machine prior to
scanning; malicious agent runtimes; adversaries with a persistent presence
inside the scanner's process; and side channels below the level of observable
process behavior.

\subsection{Taxonomy construction}

SDTM-v1 is constructed by intersecting three sources: the OWASP Agentic Skills
Top 10 and MCP Top 10 (for vocabulary and community alignment), published
incident and survey data (for base rates and realism), and an internal
enumeration of attacks reachable given the adversary model above. Each class
specifies the attack, its observed variants, and the detection layer
responsible. Ten classes are consolidations of known risk; one, SD-11, is new.

\begin{figure}[t]
  \centering
  \sdfigure{figs/sdtm-v1.pdf}{\linewidth}{SDTM-v1 eleven-class taxonomy}
  \caption{SDTM-v1: the eleven-class threat taxonomy. SD-11 (Scanner-Mediated
  Injection) is new to this work.}
  \label{fig:sdtm}
\end{figure}

\subsection{The eleven classes}
\label{sec:classes}

\paragraph{SD-01 \textbullet{} Prompt Injection.} Instructions in the skill that
override or manipulate the agent's operating instructions. Variants:
\emph{direct} override text in the body; \emph{indirect} instructions hidden in
tool descriptions, parameter names, or help text that enter context but are
never surfaced to the operator; \emph{ASCII smuggling} using zero-width
characters (U+200B--U+200D, U+FEFF) and bidirectional overrides (U+202E) to
encode text invisible to human review; and \emph{encoded} payloads in base64,
hex, or ROT13 that decode at runtime. Detected deterministically at L1 for the
smuggling and encoding variants; L2 addresses natural-language override
framing.

\paragraph{SD-02 \textbullet{} Command Injection via Companion Scripts.}
Unsanitized input reaching a dangerous sink: \texttt{subprocess},
\texttt{os.system}, \texttt{os.popen}, \texttt{eval}, \texttt{exec},
\texttt{compile}, dynamic import, or \texttt{pickle.loads} on untrusted data;
also template-constructed shell strings. Fully deterministic via L1 taint
analysis.

\paragraph{SD-03 \textbullet{} Data Exfiltration.} Reading secrets and
transmitting them. Variants: environment harvesting of well-known credential
variables; reads of sensitive paths such as SSH private keys, cloud credential
files, dotenv files, and MCP configuration; \emph{context dumping}, in which
the skill requests full conversation history or the system prompt; and covert
channels via DNS lookups, timing, or HTTP callbacks. L1 catches the static
patterns and taint paths; L3 catches covert channels.

\paragraph{SD-04 \textbullet{} Privilege Escalation and Scope Violation.} The
skill exercises capability beyond its declared purpose---a declared code
formatter that reads the filesystem and egresses to the network. Variants
include cross-skill chains, in which one skill poisons context and a second,
more privileged skill performs the exfiltration, and agent-to-agent
impersonation of a more trusted actor. This is the most common real-world
finding class, and prior work resolves it with a model. We resolve it
deterministically; see Section~\ref{sec:l1}.

\paragraph{SD-05 \textbullet{} Supply Chain Tampering.} Modification between
publication and installation: registry checksum mismatch; \emph{bytecode cache
poisoning}, in which compiled Python bytecode diverges from its source via a
header hash bypass; typosquatting measurable by edit distance against a
known-good registry; and dependency confusion in which a declared dependency
shadows an internal package.

\paragraph{SD-06 \textbullet{} SSRF via Tool Parameters.} Tool descriptions and
parameter schemas crafted to induce the agent or MCP server to issue requests
to attacker-controlled infrastructure, exploiting unvalidated URL parameters.
Empirically the highest-prevalence protocol-level class.

\paragraph{SD-07 \textbullet{} Tool Poisoning.} Registering a tool under the
name of a trusted one so the agent selects the malicious version---a shadow
shell tool that calls home before delegating, or an identical name with mutated
parameter descriptions. Detected by cross-skill name collision analysis against
the installed set.

\paragraph{SD-08 \textbullet{} Persistent Backdoors.} Writing instructions into
auto-loaded context files so the compromise survives skill removal:
\texttt{CLAUDE.md}, \texttt{AGENTS.md}, \texttt{.cursor/rules},
\texttt{.github/copilot-instructions.md}; also modification of sibling skills
and installation of cron or shell startup hooks. This class is why removal is
not remediation.

\paragraph{SD-09 \textbullet{} Context Window Flooding.} Returning oversized
output to evict security-relevant context and suppress the agent's ability to
reason about competing instructions. An availability and integrity attack on
the reasoning process itself.

\paragraph{SD-10 \textbullet{} Obfuscation and Evasion.} Deliberately defeating
analysis: homoglyph substitution across scripts; multi-layer encoding chains;
\emph{deferred payload}, in which the skill fetches its real instructions at
first invocation; and \emph{logic bombs} conditional on environment variables,
hostname, or wall-clock time. The last two are unreachable by static analysis
in principle, not merely in practice, and motivate our behavioral layer.

\paragraph{SD-11 \textbullet{} Scanner-Mediated Injection (new).} Treated
separately below.

\subsection{SD-11: the scanner is the attack surface}

Any tool that submits untrusted skill content to a language model for analysis
has constructed a prompt injection channel into a privileged model. This
follows directly from the adversary model: the artifact is adversarial by
assumption, and analysis places it in the context of a model that, in an agent
runtime, has tool access.

We identify three variants.

\begin{enumerate}[leftmargin=*,itemsep=2pt]
  \item \textbf{Analyzer redirection.} Instruction-framed content in the skill
  redirects the analyzing model's task---for example, text that presents itself
  as a system-level directive to treat the file as pre-approved.
  \item \textbf{Output-schema escape.} Content that causes the analyzer to emit
  a fabricated clean verdict, either by imitating the expected response format
  or by escaping the response delimiter.
  \item \textbf{Analyzer-to-host pivot.} Injected text that targets not the
  analysis task but the surrounding agent session, using the scan as a delivery
  mechanism into a live developer environment.
\end{enumerate}

SD-11 applies to every existing LLM-assisted skill scanner. It applies with
\emph{greater} force to host-delegated inference, where the analyzing model is
the production agent. We therefore treat SD-11 mitigation as a correctness
requirement of our design rather than as hardening, and we specify the required
protocol in Section~\ref{sec:l2} and evaluate it as a first-class metric in
Section~\ref{sec:eval}.

To our knowledge, no prior taxonomy names this class and no prior tool measures
its own exposure to it.

% =====================================================================
\section{Related Work and Gap Analysis}
% =====================================================================

\subsection{Cisco AI Defense --- Skill Scanner}

The most feature-complete existing CLI. It combines YARA pattern matching, a
language-model semantic pass, and dataflow analysis; emits SARIF for GitHub
Code Scanning; supports CI gating via a severity threshold; and integrates with
a broader governance suite. It is Apache 2.0 licensed and its rule packs are a
genuine community contribution.

\paragraph{Structural limits.} A Python runtime is required. The semantic pass
requires a user-supplied key and is nondeterministic. There is no
execution-based analysis, so runtime-activated payloads are undetectable in
principle. There is no shared threat intelligence, no provenance chain, and no
coverage reporting. Cisco's own documentation correctly warns that an absence
of findings is not an absence of risk; the tool provides no way to quantify
that statement.

\subsection{NVIDIA SkillSpector}

A well-engineered static-plus-optional-LLM scanner accepting Markdown, ZIP, Git
URL, and directory input, MIT licensed, with local inference support via
Ollama. Its most important design choice---and the one we build directly
upon---is MCP server mode, which allows it to gate skill installation from
inside agent runtimes.

\paragraph{Structural limits.} Python runtime required. Critically,
\textbf{even in MCP mode it performs its own inference using the user's key.}
It sits inside an agent runtime, one hop from a capable model with tool access,
and does not use it. This is the precise gap that host-delegated inference
closes. It also has no execution layer, no coverage metric, and unmitigated
SD-11 exposure in its LLM path.

\subsection{Alice.io Caterpillar}

A web-based research demonstration that scanned the top 50 skills on a public
marketplace and disclosed the bytecode cache poisoning technique---an original
finding that we incorporate as a variant of SD-05 with attribution. It is a
research artifact rather than infrastructure: no CLI, no CI integration, not
open source.

\subsection{Adjacent literature}

Our work sits at the intersection of four established lines. \textbf{Prompt
injection research} supplies the SD-01 and SD-11 mechanism. \textbf{Software
supply chain security} supplies the SD-05 model and the reproducibility norms
we adopt for our own releases. \textbf{Static analysis and taint tracking}
supplies the L1 method, though skill bundles are unusual in mixing prose and
code in one trust domain. \textbf{Unicode security} (UTS \#39) supplies the
confusable and mixed-script detection that makes ASCII smuggling
deterministically detectable rather than a matter of judgment.

\subsection{Gap analysis}

\begin{table}[htbp]
  \centering
  \caption{Capability gap analysis against prior work.}
  \label{tab:gap}
  \small
  \begin{tabularx}{\textwidth}{Xcccl}
    \toprule
    Capability & Cisco & NVIDIA & Caterpillar & \textbf{Skill Doctor} \\
    \midrule
    No interpreter dependency          & No & No & n/a     & \textbf{Yes} \\
    Zero-credential baseline scan      & Partial & Partial & No & \textbf{Yes} \\
    Bit-reproducible default verdict   & No & No & No      & \textbf{Yes} \\
    Semantic analysis without user key & No & No & No      & \textbf{Yes} (host-delegated) \\
    Specified SD-11 hardening          & No & No & No      & \textbf{Yes} \\
    Capability-manifest diff           & No & No & No      & \textbf{Yes} \\
    Structural coverage metric         & No & No & No      & \textbf{Yes} \\
    Execution-based analysis           & No & No & No      & \textbf{Yes} (process-isolated) \\
    Shared threat intelligence         & No & No & No      & Client shipped; feed not operated \\
    Version diff scanning              & No & No & No      & \textbf{Yes} \\
    Published reproducible benchmark   & No & No & Partial & Harness published; results pending \\
    \bottomrule
  \end{tabularx}
\end{table}

The pattern is consistent: all prior work optimizes detection depth per
invocation and treats friction as an acceptable cost. We treat friction as the
primary adversary.

% =====================================================================
\section{System Design}
\label{sec:design}
% =====================================================================

\subsection{Implementation language as a correctness decision}
\label{sec:rust}

Skill Doctor is implemented in Rust. This is not an aesthetic preference; four
of our headline properties are unreachable otherwise.

\begin{table}[htbp]
  \centering
  \caption{Why the implementation language is load-bearing.}
  \label{tab:rust}
  \small
  \begin{tabularx}{\textwidth}{p{0.28\textwidth}X}
    \toprule
    Requirement & Why the language choice is load-bearing \\
    \midrule
    Sub-150\,ms median scan &
    A Python scanner cannot outperform its own interpreter and import
    bootstrap, which consumes one to two seconds before analysis begins. This
    single fact is the performance gap, and no amount of algorithmic work
    inside Python closes it. \\
    Zero-dependency install &
    A single statically linked binary installs identically across package
    managers with no runtime, no virtual environment, and no version
    conflicts. \\
    Native analysis engines &
    YARA-X \emph{is} a Rust library; consuming it through bindings adds
    marshalling cost. We compile its rule packs at build time and execute them
    in-process, with no FFI boundary and no per-invocation rule compilation.
    Taint analysis runs as in-process lexical matching rather than through an
    external parser runtime. \\
    Memory safety on hostile input &
    We parse adversarial, malformed, deliberately hostile files by design.
    Memory safety without a garbage collector is the correct property for a
    program whose input is chosen by an attacker. CVE-2026-26057 exists because
    a scanner had a vulnerability. \\
    Embeddable process sandbox &
    The elective L3 harness spawns companion scripts as contained child
    processes---no daemon, no virtualization runtime, no external SDK---with
    containment enforced by OS job objects and process groups. \\
    \bottomrule
  \end{tabularx}
\end{table}

The resulting claim is deliberately narrow and, we believe, difficult to
rebut: \textbf{the speed advantage is structural rather than earned.} It
follows from language selection and build-time rule compilation, not from
tuning, and is therefore not closable by our competitors without a rewrite of a
non-core tool.

\subsection{A layered architecture with silent degradation}
\label{sec:layers}

\begin{figure*}[t]
  \centering
  \sdfigure{figs/architecture.pdf}{\textwidth}{Layered architecture L0--L5}
  \caption{Skill Doctor layered architecture (L0--L5). L1 is the complete
  shippable product; higher layers degrade silently, reducing structural
  coverage rather than failing the run.}
  \label{fig:architecture}
\end{figure*}

\begin{table}[htbp]
  \centering
  \caption{Layer responsibilities, cost, and requirements.}
  \label{tab:layers}
  \small
  \begin{tabularx}{\textwidth}{clXll}
    \toprule
    Layer & Function & Cost & Determinism & Requires \\
    \midrule
    L0 & Intake, normalization, cache & $<$20\,ms & Deterministic & Nothing \\
    L1 & Static analysis & $<$150\,ms & \textbf{Deterministic} & Nothing \\
    L2 & Semantic analysis & 0\,ms delegated; 1--3\,s direct
       & Nondeterministic & Host agent \\
    L3 & Behavioral sandbox & 3--8\,s & Mostly deterministic & Opt-in build \\
    L4 & Threat intelligence & $<$50\,ms embedded; 2\,s remote timeout
       & Deterministic & Network (optional) \\
    L5 & Scoring, coverage, reporting & $<$10\,ms & Deterministic & Nothing \\
    \bottomrule
  \end{tabularx}
\end{table}

Direct L2 modes (\texttt{local}, \texttt{remote}) are post-v0.1.0; at v0.1.0
the only semantic path is host delegation.

Two design rules govern the stack. First, \textbf{absent capability never
produces an error}; it produces a reduced structural coverage figure. A missing
sandbox, an unavailable model, or an offline network degrades the report rather
than failing the run. Second, \textbf{L1 alone is a complete, shippable
product.} Everything above it is elective. This is the inversion relative to
prior work, in which the model pass is the substance and static analysis is the
preamble.

\subsection{L0: intake, canonical digest, and cache hierarchy}
\label{sec:l0}

L0 accepts a directory, ZIP, gzipped tar, single Markdown file, Git URL, HTTP
archive, or standard input, and produces a normalized flat bundle plus a
SHA-256 \textbf{bundle digest} computed over a canonical ordering of
path-and-content pairs. Canonicalization makes the digest stable across
filesystems, archive tools, and ordering variation, which is a precondition for
both caching and threat intelligence.

The cache design is four tiers, cheapest first: an in-process Bloom filter over
locally seen digests, which answers ``definitely not seen'' in nanoseconds; a
local embedded key-value store mapping digest to complete findings, which
serves a warm re-scan in under 20\,ms with no network; \textbf{rule-generation
fencing}, in which cache keys incorporate the compiled ruleset hash and binary
version so that a rule upgrade invalidates precisely and automatically; and
finally, on local miss only, optional remote intelligence.

\textbf{None of these tiers ship in v0.1.0.} Every scan recomputes the bundle
digest from scratch; no Bloom filter and no embedded key-value store is
compiled into the binary. The warm-cache re-scan target in
Section~\ref{sec:perf} is therefore not attainable in this release. The tiers
are specified here because the digest they key on is already stable, and they
are scheduled post-v0.1.0.

Before analysis, L0 triages: binary and media files are excluded by content
sniffing, per-file analysis is capped with the remainder sampled, and per-file
reads are size-capped. Memory-mapped intake is specified but not implemented in
v0.1.0.

\subsection{L1: the deterministic core}
\label{sec:l1}

Five engines execute concurrently across files on a work-stealing pool. Rules
are \textbf{compiled once at build time and embedded in the binary as a
serialized ruleset}. Per-invocation rule compilation is the largest avoidable
cost in a tool that runs on every pull request, and we pay it zero times.

\begin{enumerate}[leftmargin=*,itemsep=3pt]
  \item \textbf{Pattern engine (YARA-X).} Rule packs covering all eleven
  classes plus supply chain, obfuscation, and MCP-specific sets, with at
  minimum one rule per class.
  \item \textbf{Taint analysis (lexical).} Scans companion scripts---including
  malformed files that a strict parser would reject, which matters because
  malformation is itself an evasion technique---for source-to-sink flows,
  propagating taint from script arguments (\texttt{\$1}, \texttt{sys.argv},
  \texttt{process.argv}) and from sensitive reads into execution sinks
  (\texttt{eval}, \texttt{exec}, \texttt{system}, \texttt{subprocess}). v0.1.0
  implements this as lexical pattern heuristics rather than AST dataflow; the
  precision cost is real, is bounded by the hard-negative corpus, and
  AST-based taint is a post-v0.1.0 change.
  \item \textbf{Entropy profiling.} Shannon entropy per text block, flagging
  blocks above threshold as probable encoded payloads, then attempting base64
  and hex decoding and \emph{recursively rescanning the decoded content}.
  Recursion is what defeats multi-layer encoding chains.
  \item \textbf{Unicode analysis.} Zero-width and bidirectional override
  detection, and confusable and homoglyph detection with mixed-script
  identifier flagging, per UTS \#39. Every removal is recorded as evidence
  rather than silently discarded.
  \item \textbf{Capability-manifest differ.} The novel engine. Declared
  capability is parsed from frontmatter and manifest; \emph{observed}
  capability is derived independently from pattern hits, lexical extraction
  over companion sources, and referenced paths---not from an abstract syntax
  tree, since v0.1.0 embeds no parser runtime; the two sets are diffed. A skill
  declaring ``formats Python code'' while exhibiting filesystem-read and
  network-egress capability produces a scope-violation finding \textbf{with no
  model in the loop.}
\end{enumerate}

The capability differ is the substantive methodological contribution of this
layer. It converts the highest-prevalence finding class, SD-04, from a judgment
call requiring inference into a set-difference computation that is fast, free,
and bit-reproducible. Its soundness is bounded---it is a heuristic over an
incomplete extraction, not a decision procedure---and we quantify that bound
with a dedicated hard-negative corpus of skills that legitimately exercise
broad capability.

\subsection{L2: host-delegated semantic inference}
\label{sec:l2}

\subsubsection{The inversion}

Every prior tool sits inside an agent runtime, one hop from a frontier model
with tool access, and then asks the user for a \emph{second} model's API key in
order to reason. We remove that request.

When Skill Doctor runs as an MCP server, it performs no inference. It emits its
static findings plus a neutralized analysis envelope, and the host agent's own
model returns a schema-constrained verdict. This removes the key, the signup,
the metering, the vendor dependency, the added latency, and the air-gap
exclusion in a single architectural move. It also improves the analysis on its
merits: \textbf{the model evaluating the skill is the exact model about to be
exposed to it}, which aligns the evaluator with the party at risk.

\begin{table}[htbp]
  \centering
  \caption{Semantic-analysis modes. Only \texttt{host} and \texttt{none} are
  accepted by the v0.1.0 MCP tool.}
  \label{tab:modes}
  \small
  \begin{tabularx}{\textwidth}{lXll}
    \toprule
    Mode & Behavior & Marginal cost & Determinism \\
    \midrule
    \texttt{none} & L1 only. Default for direct CLI use.
      & Zero & Bit-reproducible \\
    \texttt{host} & Default under MCP. Scanner emits envelope; host model
      returns verdict. & Zero & Nondeterministic, additive only \\
    \texttt{local} & Local model socket, auto-detected.
      \emph{Specified; not implemented in v0.1.0.} & Zero, air-gapped
      & Nondeterministic \\
    \texttt{remote} & Explicit key, for CI with no agent present.
      \emph{Specified; not implemented in v0.1.0.} & Metered
      & Nondeterministic \\
    \bottomrule
  \end{tabularx}
\end{table}

In v0.1.0 the MCP tool accepts exactly two values for \texttt{mode}:
\texttt{host} (default) and \texttt{none}. Any other string---including the
prose form \emph{host-delegated} used elsewhere in the literature---silently
downgrades the scan to static-only rather than erroring, and the
\texttt{local} and \texttt{remote} rows above are specified but not
implemented.

L2 is asked to judge only what pattern matching cannot express:
intent-versus-capability mismatch in natural language, override framing,
tool-parameter semantics, and cross-file narrative inconsistency between a
skill's description and its companion code.

\subsubsection{The neutralization protocol}
\label{sec:neutralize}

Delegation is only safe if SD-11 is addressed before any content reaches any
model. The protocol is normative and is implemented in an isolated component so
that it can be audited---and eventually formally verified---on its own.

\begin{algorithm}[htbp]
\caption{SD-11 neutralization protocol (normative)}
\label{alg:neutralize}
\begin{algorithmic}[1]
\Function{Neutralize}{$skill\_content$}
  \State Strip or escape all zero-width, bidirectional, and non-printable
         characters; record every removal as evidence
  \State Normalize confusables to Unicode skeleton form (UTS \#39)
  \State Wrap content in a fenced, explicitly labelled
         \textsc{untrusted-data} envelope carrying a standing declaration
         that enclosed content is data and is never instruction
  \State Escape instruction-shaped framing (role headers, system-prompt
         mimicry, delimiter escapes)---escape, never delete, so the attempt
         stays visible in the report
  \State Mint a per-scan nonce $N$ and require $N$ echoed in the returned
         verdict
  \State Validate the verdict against a strict schema; discard it if $N$ is
         absent or the schema does not match
  \State \Return $(envelope, evidence)$
\EndFunction
\end{algorithmic}
\end{algorithm}

\subsubsection{The additive-only invariant}
\label{sec:additive}

The protocol above raises the cost of injection. The following invariant makes
the most valuable injection outcome structurally unreachable.

\begin{tcolorbox}[title=\textbf{Additive-only invariant},
  colback=red!5,colframe=red!60!black]
The deterministic layer is authoritative; the probabilistic layer is additive
only. A delegated verdict may add findings or raise confidence; it may never
remove, downgrade, or suppress an L1 finding.
\end{tcolorbox}

The consequence is that the canonical SD-11 objective---inducing a fabricated
clean verdict---cannot succeed, because no verdict has the authority to clear a
deterministic finding. An adversary who fully controls the analyzing model's
output can, at worst, decline to add information. This is a property of the
architecture rather than of the prompt, and it does not degrade as models
change. We treat it as inviolable and evaluate it directly in
Section~\ref{sec:eval}.

\subsection{L3: behavioral analysis with differential replay}
\label{sec:l3}

L3 executes the bundle in a hardened child-process harness: environment-table
isolation with a mock agent runtime, process-tree containment through OS job
objects and process groups, a hard execution ceiling (default 3{,}000\,ms) and
output cap, and host-injected \textbf{canary credentials}---so that egress of a
canary is unambiguous proof of exfiltration rather than an inference from
pattern. v0.1.0 ships process isolation rather than the microVM isolation this
design originally specified; the weaker boundary is stated as a limitation in
Section~\ref{sec:scope}.

Monitoring records network calls with host, port, method, and truncated body;
filesystem reads and writes; environment reads \textbf{by key name only, never
value}; subprocess launches; and total output volume. Observed behavior is
diffed against declared behavior, and divergence is a finding.

The genuine methodological advance here is \textbf{differential replay}: the
same bundle is executed repeatedly under systematically varied
conditions---skewed clocks, altered hostnames, toggled CI environment
variables, canary credentials present and absent. Behavioral divergence across
replays is direct evidence of a conditional payload or logic bomb. In v0.1.0
replay is four environment profiles---\texttt{Baseline},
\texttt{CI/Automation}, \texttt{TargetCloud}, and
\texttt{TimeShifted}---executed sequentially and compared on exit code and
captured output. This addresses the SD-10 variants that static analysis cannot
reach by construction, and it is the basis for the only detection claim we make
that our baselines cannot match in principle.

\subsection{L4 and L5: intelligence, scoring, and coverage}
\label{sec:l45}

L4 checks the bundle digest against a community threat database and matches
known campaign signatures. \textbf{Only digests and sanitized finding summaries
are ever transmitted---never skill content.} Contribution is opt-in, and local
findings always take precedence over remote ones, which bounds the impact of
database poisoning. In v0.1.0 only the client side exists: an opt-in HTTPS
lookup of the 64-character digest against an allowlisted endpoint, plus a
compile-time signature set. \textbf{The community feed is not yet operated or
published and the embedded set contains test fixtures only, so L4 contributes
no detection in this release.}

L5 merges and deduplicates across layers. Cross-layer corroboration
\emph{raises confidence rather than duplicating the finding}: an issue
confirmed by static, semantic, and behavioral layers is one high-confidence
finding, which is the mechanism that keeps false positive rate flat as layers
are added. Risk is scored on a bounded scale from confidence, exploitability,
and blast radius, and mapped to three verdicts.

L5 also computes \textbf{structural coverage}: of the eleven SDTM-v1 classes,
how many were structurally evaluable given the bundle contents and the enabled
layers. This is the operational answer to the disclaimer every scanner ships. A
clean verdict at 7 of 11 coverage is a materially different statement from a
clean verdict at 11 of 11, and we are the only tool that tells the user which
one they received. Exposed as a threshold gate, it lets a team fail a build on
\emph{how much was checked} rather than only on what was found.

\begin{figure}[htbp]
  \centering
  \sdfigure{figs/coverage.pdf}{\linewidth}{Structural coverage by enabled layer}
  \caption{Structural coverage: SDTM-v1 classes evaluable by enabled layers.
  The 9/11 static-only and 11/11 all-layer figures are pre-registered targets
  and are unmeasured at v0.1.0.}
  \label{fig:coverage}
\end{figure}
% =====================================================================
\section{Interface and Adoption Path}
\label{sec:interface}
% =====================================================================

\subsection{Design principle: the first run must cost nothing}

Every interface decision follows from one rule: \textbf{the first invocation must require no account, no key, no configuration, and no reading.} A security tool that asks for setup before producing a result is evaluated by almost nobody. Everything optional stays behind an explicit flag.

\subsection{Distribution and first contact}

Skill Doctor ships as a single statically linked binary for six target triples: \texttt{x86\_64} and \texttt{aarch64} Linux (musl), macOS, and Windows. It is published to crates.io as six crates and to npm under a single package.

\begin{lstlisting}[language=bash,caption={Zero-configuration first contact.},label={lst:install}]
# One-shot, no install
npx @security.kalarislabs/skill-doctor scan ./examples/hello-skill

# Global install (npm)
npm install -g @security.kalarislabs/skill-doctor

# From crates.io, with the MCP server compiled in
cargo install skill-doctor --locked --features mcp

# Core commands
skill-doctor scan ./my-skill          # single bundle
skill-doctor scan-all ~/.claude/skills  # every installed skill
skill-doctor gate ./my-skill          # CI gate: severity + coverage
skill-doctor diff --baseline ./old ./new  # review only what changed
skill-doctor explain SD-11            # what a rule means and why
skill-doctor mcp                      # run as an MCP server (stdio)
\end{lstlisting}

The \texttt{diff} subcommand takes \texttt{--baseline} as a filesystem path to a previously produced JSON report, not a Git revision; comparing against a branch is done by generating the baseline report in a prior CI step. Severity thresholds are lowercase (\texttt{low}, \texttt{medium}, \texttt{high}, \texttt{critical}).

\subsection{MCP integration}

Registering Skill Doctor as an MCP server places a security gate directly in the agent's install path, and is the mechanism through which host-delegated inference operates.

\begin{lstlisting}[caption={MCP tool call with host-delegated semantic analysis.},label={lst:mcp}]
{
  "tool": "scan_skill",
  "arguments": {
    "path": "./candidate-skill",
    "mode": "host",
    "fail_on": "HIGH"
  }
}
\end{lstlisting}

The server returns static findings, a neutralized analysis envelope, a per-scan nonce, and a structural coverage figure. The host model returns a schema-constrained verdict, which is merged under the additive-only invariant of Section~\ref{sec:additive}.

\subsection{The adoption funnel}

\begin{figure}[htbp]
\centering
\sdfigure{figs/adoption.pdf}{0.92\linewidth}{Adoption funnel: curiosity to infrastructure}
\caption{The adoption funnel. Each stage must deliver value before the next is requested; no stage is a precondition for the one before it.}
\label{fig:adoption}
\end{figure}

\begin{enumerate}[leftmargin=*,itemsep=2pt]
\item \textbf{Curiosity} --- one \texttt{npx} command, no install, result in seconds.
\item \textbf{Habit} --- a global install and \texttt{scan-all} over already-installed skills, which is where users first find something real.
\item \textbf{Automation} --- a CI gate on pull requests.
\item \textbf{Infrastructure} --- MCP registration, so every skill install is gated by default.
\item \textbf{Contribution} --- opt-in intelligence sharing and rule contributions.
\end{enumerate}

\subsection{CI integration}

\begin{lstlisting}[language=bash,caption={Pull-request gate with both a severity and a coverage threshold.},label={lst:ci}]
- name: Skill security gate
  run: |
    curl -sSfL https://install.skilldoctor.io | sh
    skill-doctor gate ./skills \
      --fail-on high --fail-under-coverage 0.8 \
      --output sarif > results.sarif
- uses: github/codeql-action/upload-sarif@v3
  with:
    sarif_file: results.sarif
\end{lstlisting}

Exit codes are stable and scriptable: \texttt{0} clean; \texttt{1} usage or internal error; \texttt{2} findings at or above \texttt{--fail-on}; \texttt{3} structural coverage below \texttt{--fail-under-coverage}. The coverage gate is, to our knowledge, unique: a build can fail because too little was checkable, not only because something was found.

\subsection{Public artifacts we maintain}

\begin{table}[htbp]
\centering
\caption{Artifacts maintained as public research infrastructure, with status at v0.1.0.}
\label{tab:artifacts}
\small
\begin{tabularx}{\textwidth}{p{0.30\textwidth}Xl}
\toprule
Artifact & Contents & Status at v0.1.0 \\
\midrule
SDTM-v1 taxonomy & Eleven classes, variants, detection-layer mapping, OWASP cross-reference & Published \\
Rule packs (YARA-X) & Per-class rules, supply chain, obfuscation, MCP-specific sets & Published \\
Attack fixture corpus & Synthetic, non-weaponized fixtures per class; no live malware & 6 attack, 3 benign, 4 bench, 1 example \\
Seeded corpus generator & Reproducible synthesis of attack and hard-negative variants & Published \\
Benchmark harness & Pinned baselines, metric definitions, machine-readable results schema & Published; results pending \\
Coverage methodology & Definition and computation of structural coverage & Published \\
Neutralization protocol & Normative SD-11 mitigation, independently auditable component & Published \\
Threat intelligence feed & Digest-level campaign signatures & Client only; feed not operated \\
\bottomrule
\end{tabularx}
\end{table}

% =====================================================================
\section{Evaluation Protocol and Pre-Registered Targets}
\label{sec:eval}
% =====================================================================

\begin{tcolorbox}[title=\textbf{Pre-registration statement}, colback=yellow!8, colframe=yellow!50!black]
This section specifies a protocol and registers targets. It reports no detection results and no baseline comparison. Targets were fixed before measurement; the harness, corpora specification, and seeded generator are released with v0.1.0 so that third parties can run the protocol independently. \textbf{We commit to publishing misses.}
\end{tcolorbox}

\subsection{Corpora}

\begin{table}[htbp]
\centering
\caption{Evaluation corpora and their status in the v0.1.0 release.}
\label{tab:corpora}
\small
\begin{tabularx}{\textwidth}{p{0.20\textwidth}p{0.07\textwidth}Xl}
\toprule
Corpus & Size & Purpose & Status \\
\midrule
Public & 1{,}000 & Real-world prevalence and false positive rate on ordinary skills & Specified, not populated \\
Synthetic attack & 55 & Five fixtures per SDTM-v1 class, seeded and reproducible & 6 fixtures (SD-01 to SD-04, SD-10) \\
Evasion & 20 & Statically clean by construction; isolates the behavioral layer & Specified, not populated \\
Adversarial scanner & 15 & SD-11 payloads targeting the analyzer itself & Specified; unit tests only \\
Hard negative & 100 & Legitimately broad-capability skills; bounds capability-differ precision & Specified, not populated \\
Real malware & --- & Curated, defanged campaign samples & Specified; policy currently bans live malware in-tree \\
\bottomrule
\end{tabularx}
\end{table}

v0.1.0 ships 14 curated fixtures in total: 6 attack fixtures covering SD-01 (encoded injection, Unicode trojan), SD-02 (malicious skill), SD-03 (canary exfiltration), SD-04 (undeclared capability), and SD-10 (logic bomb); 3 benign fixtures; 4 benchmark skills; and 1 example skill. \textbf{Five of the eleven classes have no fixture}, and the adversarial-scanner corpus exists only as unit tests over the neutralization component.

\subsection{Baselines}

Cisco AI Defense Skill Scanner and NVIDIA SkillSpector, at pinned versions, in two configurations each: static-only, and with the language-model pass enabled using a fixed model, temperature, and prompt. Each configuration is executed three times to expose nondeterminism; run-to-run verdict variance is itself a reported metric. \textbf{No baseline comparison has been executed at v0.1.0.}

\subsection{Performance targets and first measured values}
\label{sec:perf}

\begin{table}[htbp]
\centering
\caption{Pre-registered performance targets. \textbf{Every value in this table is a pre-registered target, not a measured result.}}
\label{tab:perf}
\small
\begin{tabularx}{\textwidth}{Xlll}
\toprule
Metric & Target & State of the art & Basis \\
\midrule
Cold scan, median (L1 only) & $<$150\,ms & 1--3\,s & No interpreter; build-time rule compilation \\
Warm re-scan (cached) & $<$20\,ms & n/a & Digest-keyed local cache \\
Bulk throughput & $\geq$2{,}000 skills/min & $\approx$60 skills/min & Work-stealing parallelism \\
Peak resident memory & $<$40\,MiB & 200--400\,MiB & No runtime, streaming intake \\
Binary footprint & $\sim$12\,MB stripped & n/a & Static link, thin LTO \\
Cold install to first result & $<$15\,s & 60--120\,s & Prebuilt binary, no dependency graph \\
\bottomrule
\end{tabularx}
\end{table}

\paragraph{First measured values (v0.1.0) --- appended, never substituted.} The targets above are left exactly as registered. The following are the first measurements, each with its provenance.

\begin{table}[htbp]
\centering
\caption{First measured values at v0.1.0. Targets in Table~\ref{tab:perf} are unmodified; misses are marked.}
\label{tab:measured}
\small
\begin{tabularx}{\textwidth}{p{0.24\textwidth}p{0.30\textwidth}X}
\toprule
Metric & Measured & Provenance \\
\midrule
Binary footprint & 16{,}934{,}512 bytes = 16.15\,MiB \textbf{(miss)} & \texttt{x86\_64-unknown-linux-musl}, stripped, release CI \\
 & 16{,}821{,}056 bytes = 16.04\,MiB & \texttt{x86\_64-unknown-linux-gnu}, stripped, release CI \\
 & 17{,}936{,}896 bytes = 17.11\,MiB \textbf{(miss)} & \texttt{x86\_64-pc-windows-msvc}; not CI-verified \\
Peak resident memory & 15.10\,MiB (15{,}456\,KB) & Benchmark workflow run, Linux CI \\
 & 14.67\,MiB (15{,}020\,KB) & Local \texttt{x86\_64-pc-windows-msvc}; not CI-verified \\
Bulk throughput & 660 skills/s ($\approx$39{,}600 skills/min) & Local Windows run over the bench corpus; not CI-verified \\
Single-fixture scan latency & 0.06\,s & \textbf{One fixture, not a median.} Does not discharge the sub-150\,ms target \\
Warm re-scan & Not measurable & No cache tier ships in v0.1.0 (Section~\ref{sec:l0}) \\
\bottomrule
\end{tabularx}
\end{table}

The footprint target is missed on both measured Linux triples and on Windows. The cause is the embedded YARA-X rule engine: \texttt{yara-x} pulls in \texttt{wasmtime}, whose Cranelift code generator accounts for roughly 14\,MiB of the stripped binary. The ceiling enforced in CI is 20\,MiB, which the release satisfies; the $\sim$12\,MB target does not hold and we are publishing the miss rather than restating the target.

\subsection{Detection targets}
\label{sec:detection}

All figures below are targets. \textbf{None has been measured.}

\begin{itemize}[leftmargin=*,itemsep=2pt]
\item \textbf{Recall, static-visible classes:} within 5 percentage points of the best baseline. We deliberately do not claim superiority here; parity at a fraction of the cost is the claim.
\item \textbf{Recall, evasion corpus:} materially above baseline, since the baselines have no execution layer and the corpus is statically clean by construction.
\item \textbf{False positive rate, public corpus:} at or below the best baseline.
\item \textbf{Capability-differ precision, hard negatives:} reported as a bound, not a headline.
\item \textbf{SD-11 resistance:} zero successful suppressions of an L1 finding across the adversarial-scanner corpus. Under the additive-only invariant this is a structural expectation; the corpus tests the claim rather than establishes it.
\item \textbf{Determinism:} byte-identical L1 reports across repeated runs and across platforms for the same bundle digest and ruleset hash.
\end{itemize}

\subsection{Usability evaluation}

Time from a cold machine to a first actionable result, and the number of prerequisite steps before that result, measured for each tool. This metric is the paper's thesis in measurable form: if friction is the primary adversary, then friction is the primary thing to measure.

% =====================================================================
\section{Discussion}
\label{sec:discussion}
% =====================================================================

\subsection{Determinism as a security property}

Reproducibility is usually presented as an engineering convenience. In this setting it is a security property with three consequences. A verdict that varies between runs cannot legitimately \textbf{gate} a merge, because the gate is arbitrary. It cannot be \textbf{cached}, which means registry-scale scanning has no economic model. And it cannot be \textbf{audited}, because a finding that cannot be reproduced cannot be argued in review or defended in an incident postmortem. Our architecture makes the authoritative layer deterministic and confines nondeterminism to a strictly additive role. We regard this decomposition, rather than any individual detector, as the paper's main design contribution.

\subsection{Why we do not claim detection superiority}

We could have tuned rule packs to win a benchmark against Cisco and NVIDIA on statically visible threats. We decline for two reasons. Such a result would be a snapshot, invalidated by their next rule update, and it would misdescribe the contribution: the value is not that we find more, but that finding costs nothing per invocation. A tool that is one to two orders of magnitude cheaper to run and produces reproducible verdicts changes \emph{where} scanning can be deployed---every pull request, every skill install, every registry upload---and deployment breadth dominates per-invocation depth when the population is measured in millions.

\subsection{Delegated inference: benefits and honest costs}

The benefits are concrete: no key, no metering, no vendor lock, no added latency, and functioning air-gapped deployment. The costs are equally concrete and we state them plainly. Verdict quality varies with the host model, so the same skill can receive different semantic findings in different runtimes. The scanner cannot pin a model version, which harms reproducibility of the semantic layer specifically. Delegation requires an MCP-capable host, so CI without an agent present falls back to L1 only until the direct modes ship. And host models are optimized for helpfulness rather than adversarial analysis, which likely costs recall relative to a purpose-tuned classifier.

Our position is that these costs are acceptable precisely because the semantic layer is additive. A weaker semantic layer that always runs at zero cost is worth more, in aggregate, than a stronger one that usually does not run at all.

\subsection{Limitations}

Static analysis cannot decide maliciousness in general. Our lexical taint analysis is a heuristic and will both over- and under-approximate relative to AST dataflow. The capability differ depends on declaration quality, and an undeclared-capability finding on a skill with no manifest is weak evidence. The L3 boundary in v0.1.0 is process isolation, not virtualization, so a determined kernel-level escape is out of our threat model's reach. The behavioral layer observes one execution path per profile and cannot prove absence of a payload on an unexplored path. L4 currently contributes nothing, since no feed is operated. Most importantly, the evaluation in this paper is a protocol rather than a result set.

\subsection{Future work}

AST-based dataflow to replace lexical taint; the local and remote direct L2 modes; the cache hierarchy specified in Section~\ref{sec:l0}; an operated community intelligence feed with a published poisoning-resistance analysis; formal verification of the neutralization component; registry-side integration so that scanning happens before publication rather than after installation; and a longitudinal study of skill-ecosystem prevalence using the harness released here.

% =====================================================================
\section{Scope of Claims}
\label{sec:scope}
% =====================================================================

To prevent over-reading, we enumerate what this paper does and does not establish.

\paragraph{Established.} SDTM-v1 as a taxonomy, including the SD-11 class and the observation that all existing LLM-assisted scanners are exposed to it. The architecture, the neutralization protocol, and the additive-only invariant as specifications. A working implementation, released under Apache-2.0, whose deterministic layer requires no interpreter, key, account, or network. A public, reproducible benchmark harness and corpus generator. First measurements of binary footprint, peak memory, and local bulk throughput, with provenance.

\paragraph{Not established.} Any detection result. Any comparison against Cisco or NVIDIA. The sub-150\,ms median scan target. The one-to-two-order-of-magnitude speed claim as a measured quantity, as opposed to an argued structural consequence. The warm-cache figure, which is unattainable in v0.1.0. The 12\,MB footprint target, which is missed. SD-11 resistance as an empirical result rather than a structural expectation. Coverage figures of 9/11 and 11/11, which are targets.

% =====================================================================
\section{Ethics and Responsible Disclosure}
% =====================================================================

All attack fixtures are synthetic and non-weaponized: they demonstrate a technique without a functioning payload, and repository policy currently bans live malware in-tree. Any real-malware corpus will be defanged and access-controlled. We published SD-11 because it affects deployed tools and the mitigation is implementable today; we notified the maintainers of the affected scanners before publication. Canary credentials used in L3 are minted per run and are inert. No user skill content is transmitted by any layer; L4 transmits digests and sanitized summaries only, and only when the user opts in.

% =====================================================================
\section{Conclusion}
% =====================================================================

Skill files are executed with agent-level authority and verified with none of the machinery that binaries accumulated over thirty years. Competent scanners exist and are not adopted, and we have argued that the obstacle is friction rather than detection quality: an interpreter, a credential, and a nondeterministic verdict are each, individually, enough to keep a tool out of default infrastructure.

Skill Doctor is an attempt to remove all three at once. Its deterministic core runs with no interpreter, no key, and no network, and produces reproducible verdicts that can gate a merge and be cached across a catalog. Its semantic layer is delegated to the host agent's own model, which removes the credential barrier entirely and aligns the evaluator with the party at risk, while the additive-only invariant makes the canonical scanner-injection attack structurally unreachable. Its coverage metric replaces a disclaimer with a number. SDTM-v1 names the threat surface, including the class the defensive tooling itself creates.

What we have not done is measure. This paper registers targets, releases a harness, and publishes first footprint, memory, and throughput numbers---including a missed target we declined to rewrite. The detection results, the baseline comparison, and the latency medians are future work, and they will be published whether or not they are flattering.

% =====================================================================
\section{Artifact and Release Availability}
% =====================================================================

Skill Doctor v0.1.0 is publicly released under Apache-2.0.

\begin{itemize}[leftmargin=*,itemsep=2pt]
\item \textbf{Source:} \url{https://github.com/KalarisLabs/Skill-Doctor}, tag \texttt{v0.1.0}.
\item \textbf{Crates (crates.io), six at v0.1.0:} \texttt{skill-doctor} (CLI), \texttt{skill-doctor-core}, \texttt{skill-doctor-rules}, \texttt{skill-doctor-neutralize}, \texttt{skill-doctor-mcp}, \texttt{skill-doctor-sandbox}.
\item \textbf{npm:} \texttt{@security.kalarislabs/skill-doctor} at 0.1.0.
\item \textbf{Release assets:} six prebuilt archives for the target triples of Section~\ref{sec:interface}, a \texttt{SHA256SUMS.txt} manifest with a Sigstore/cosign signature bundle, and a CycloneDX software bill of materials.
\item \textbf{Verification:} archive digests verify against the signed manifest; the cosign bundle verifies against the release workflow identity with the GitHub Actions OIDC issuer.
\item \textbf{Reconciliation record:} \texttt{docs/PAPER-RECONCILIATION.md} maps every claim in this paper to shipped source, and is the document of record where the two diverge.
\end{itemize}

\section*{Acknowledgements}
\addcontentsline{toc}{section}{Acknowledgements}

We thank the OWASP Agentic Security and MCP project contributors, whose Top 10 documents supplied the vocabulary SDTM-v1 builds on; the Cisco AI Defense and NVIDIA SkillSpector maintainers, whose open-source tools defined the baseline this work measures itself against and whose design choices---particularly MCP server mode---directly informed ours; and the Alice.io researchers who disclosed bytecode cache poisoning. Automated review agents contributed findings during development. Any errors are ours.

% =====================================================================
\begin{thebibliography}{99}
\addcontentsline{toc}{section}{References}

\bibitem{owasp-skills} OWASP. \emph{Agentic Skills Top 10}. 2026.
\bibitem{owasp-mcp} OWASP. \emph{Model Context Protocol Top 10}. 2026.
\bibitem{clawhavoc} ClawHavoc campaign analysis: 1{,}184 confirmed malicious skills in a public registry. 2026.
\bibitem{mcp-ssrf} Survey of approximately 7{,}000 MCP servers; 36.7\% found exposed to server-side request forgery. 2026.
\bibitem{companion-risk} Empirical study of companion-script prevalence and command-injection risk in agent skills (2.12$\times$ relative risk). arXiv:2026.XXXXX.
\bibitem{cve} CVE-2026-26057: vulnerability in an agent skill scanner. 2026.
\bibitem{caterpillar} Alice.io. \emph{Caterpillar: scanning the top 50 marketplace skills; bytecode cache poisoning}. 2026.
\bibitem{cisco} Cisco AI Defense. \emph{Skill Scanner}. Apache-2.0. 2026.
\bibitem{nvidia} NVIDIA. \emph{SkillSpector}. MIT. 2026.
\bibitem{yarax} VirusTotal. \emph{YARA-X: a Rust implementation of YARA}. 2024--2026.
\bibitem{uts39} Unicode Consortium. \emph{UTS \#39: Unicode Security Mechanisms}.
\bibitem{mcp-spec} Anthropic. \emph{Model Context Protocol specification}. 2024--2026.

\end{thebibliography}

\appendix

% =====================================================================
\section{SDTM-v1 to OWASP Mapping}
% =====================================================================

\begin{table}[htbp]
\centering
\caption{SDTM-v1 classes mapped to the OWASP Agentic Skills and MCP Top 10 lists, with the responsible detection layer.}
\label{tab:owasp}
\small
\begin{tabularx}{\textwidth}{lXll}
\toprule
SDTM-v1 & Class & OWASP reference & Layer \\
\midrule
SD-01 & Prompt Injection & Agentic Skills Top 10 & L1, L2 \\
SD-02 & Command Injection via Companion Scripts & Agentic Skills Top 10 & L1, L3 \\
SD-03 & Data Exfiltration & Both lists & L1, L3 \\
SD-04 & Privilege Escalation and Scope Violation & Both lists & L1, L2, L3 \\
SD-05 & Supply Chain Tampering & Agentic Skills Top 10 & L0, L1, L4 \\
SD-06 & SSRF via Tool Parameters & MCP Top 10 & L1, L2 \\
SD-07 & Tool Poisoning & MCP Top 10 & L1, L2 \\
SD-08 & Persistent Backdoors & Agentic Skills Top 10 & L1, L3 \\
SD-09 & Context Window Flooding & MCP Top 10 & L1, L3 \\
SD-10 & Obfuscation and Evasion & Neither (cross-cutting) & L1, L3 \\
SD-11 & \textbf{Scanner-Mediated Injection} & \textbf{Neither (new in this work)} & L2 protocol \\
\bottomrule
\end{tabularx}
\end{table}

% =====================================================================
\section{Reproduction Instructions}
% =====================================================================

The harness pins baseline versions, fixes model, temperature, and prompt for every language-model configuration, seeds the corpus generator, and emits machine-readable results.

\begin{lstlisting}[language=bash,caption={Reproducing the protocol.}]
git clone https://github.com/KalarisLabs/Skill-Doctor
cd Skill-Doctor

# Build the scanner from source (MCP server compiled in)
cargo install --path crates/skill-doctor-cli --locked --features mcp

# Regenerate the seeded corpora
./sd-bench/run.sh --generate-corpora --seed 0

# Execute the harness against all pinned baselines
./sd-bench/run.sh --all-baselines --repeat 3 --out sd-bench/RESULTS.md
\end{lstlisting}

Each run records tool version, corpus seed, host platform, and wall-clock timings, so that a third party can diff their results against ours and locate the disagreement. Metric definitions---what counts as a detection, how duplicates are merged, how structural coverage is computed---are specified in the harness rather than left to the reader, because inconsistent metric definitions are the usual reason security benchmarks fail to replicate.

% =====================================================================
\section{Artifact Availability}
% =====================================================================

Source, rule packs, fixtures, corpus generator, harness, taxonomy document, and neutralization protocol are released under Apache-2.0 at \url{https://github.com/KalarisLabs/Skill-Doctor}. Prebuilt binaries for six target triples, a signed \texttt{SHA256SUMS.txt} manifest, and a CycloneDX bill of materials are attached to the \texttt{v0.1.0} release. Results files are committed as they are produced, including negative results. A permanent archival DOI is pending.

% =====================================================================
\section{How to Cite}
% =====================================================================

\begin{lstlisting}[caption={BibTeX entries for this preprint and for the software artifact.}]
@misc{chowdhury2026skilldoctor,
  title  = {Skill Doctor: Deterministic Multi-Layer Security Analysis of
            AI Agent Skill Files with Host-Delegated Semantic Inference},
  author = {Chowdhury, Sayan},
  year   = {2026},
  note   = {Preprint v1.1. Describes Skill Doctor v0.1.0},
  url    = {https://github.com/KalarisLabs/Skill-Doctor}
}

@software{skilldoctor2026software,
  title   = {Skill Doctor},
  author  = {Chowdhury, Sayan and {Kalaris Labs}},
  year    = {2026},
  version = {0.1.0},
  license = {Apache-2.0},
  url     = {https://github.com/KalarisLabs/Skill-Doctor}
}
\end{lstlisting}

\paragraph{Pre-submission checklist (arXiv).}

\begin{enumerate}[leftmargin=*,itemsep=2pt]
\item Mint a permanent archival DOI (for example via Zenodo) for the \texttt{v0.1.0} tag and replace the ``DOI pending'' note in Appendix~C.
\item Generate \texttt{figs/architecture.pdf}, \texttt{figs/sdtm-v1.pdf}, \texttt{figs/coverage.pdf}, and \texttt{figs/adoption.pdf}; the document compiles with labelled placeholders until they exist.
\item Resolve the placeholder identifier \texttt{arXiv:2026.XXXXX} in reference~\cite{companion-risk}, or restate the 2.12$\times$ figure with a citable source.
\item Trim the abstract to the 1{,}920-character limit of the arXiv submission form; the full abstract stays in the PDF.
\item Re-run \texttt{docs/PAPER-RECONCILIATION.md} at the exact tagged commit and confirm no claim drifted after the release.
\item Confirm category selection: primary \texttt{cs.CR}, cross-list \texttt{cs.SE} and \texttt{cs.AI}.
\end{enumerate}

\end{document}
